\section{Conclusion}
\label{sec:conclusion}

This work provides evidence that episodic memory is not a universal enhancement for LLM-based agents. Through a controlled study applying an identical memory system to ALFWorld and WebShop, we observe that the same retrieval mechanisms produce a ${\sim}47$pp swing in effect: up to $+19.4$pp improvement versus $-30.5$pp degradation. Our hard-negative ablation reveals that in the low-transfer environment we studied, irrelevant memories caused \emph{less} harm than semantically ``relevant'' ones, suggesting confident misdirection rather than noise injection as the failure mode — at least under these experimental conditions.

We attribute this divergence to three structural properties (task compositionality, action space regularity, and learning transfer radius) and contribute a diagnostic framework offering practitioners tentative criteria for anticipating when episodic memory may help or hurt. Because these criteria are grounded in two environments with a single base model, they should be treated as hypotheses to test rather than universal rules. Future directions include adaptive memory gating that suppresses retrieval when transfer confidence is low, validation on additional domains and models, and memory filtering strategies that assess transferability before injecting retrieved content.

\paragraph{AI Disclosure.}
AI writing assistants were used to help draft and edit portions of this manuscript. All scientific content, experimental design, and analysis were conducted by the authors.
